% CORRECTED VERSION
% Changes:
% - Fixed misuse of Lemma 2 (independence requirement) by proving a new sign‑mismatch bound using Chebyshev’s inequality.
% - Replaced the unjustified Chebyshev step with a rigorous variance recursion for the momentum term.
% - Added Lemma 3 (variance of the EMA) and Lemma 4 (sign‑mismatch probability).
% - Corrected the lower bound on the expected inner product (Step 8) and removed the invalid dimension‑free claim.
% - Made the convergence rate consistent by using a stepsize α = c/√T (still ≤ 1/L) and explicitly deriving the O(1/√T) bound.
% - Stated clearly that  min ≤ average to justify the final “minimum’’ guarantee.
% - Updated Theorem 1 accordingly and clarified all assumptions.

\documentclass{article}
\usepackage{amsmath,amssymb,amsthm}
\usepackage{algorithm}
\usepackage{algorithmic}
\usepackage{xcolor}
\newtheorem{theorem}{Theorem}
\newtheorem{lemma}{Lemma}
\newtheorem{assumption}{Assumption}
\newcommand{\sign}{\operatorname{sign}}
\newcommand{\E}{\mathbb{E}}
\newcommand{\R}{\mathbb{R}}

\begin{document}

\section*{SignSGD with Momentum}

\section*{Mathematical Formulation}
Let $f:\R^{d}\rightarrow\R$ be differentiable.  
At iteration $k\ge 0$ we have a stochastic gradient $g_{k}= \nabla f(x_{k};\xi_{k})$, where $\xi_{k}$ denotes the random data sample.  
The algorithm maintains a momentum vector $v_{k}\in\R^{d}$ and updates as follows
\begin{align}
v_{k} &= \beta\,v_{k-1} + (1-\beta)\,g_{k},\qquad v_{-1}=0,\label{eq:momentum}\\[2mm]
d_{k} &= \sign(v_{k})\in\{-1,0,+1\}^{d},\label{eq:direction}\\[2mm]
x_{k+1} &= x_{k} - \alpha_{k}\, d_{k}, \qquad \alpha_{k}>0 \text{ (stepsize).}\label{eq:update}
\end{align}
Thus only the sign of the momentum‑averaged gradient is used to move each coordinate by a (possibly time‑varying) amount $\alpha_{k}$.

\section*{Pseudocode}
\begin{algorithm}[H]
\caption{SignSGD with Momentum}
\begin{algorithmic}[1]
\STATE \textbf{Input:} base stepsize $c>0$, momentum parameter $\beta\in[0,1)$, max iter $T$.
\STATE Initialize $x_{0}\in\R^{d}$, $v_{-1}=0$.
\FOR{$k=0$ \textbf{to} $T-1$}
    \STATE Sample a minibatch $\xi_{k}$ and compute stochastic gradient $g_{k}= \nabla f(x_{k};\xi_{k})$.
    \STATE $v_{k}= \beta v_{k-1} + (1-\beta) g_{k}$ \hfill \textcolor{gray}{// \eqref{eq:momentum}}
    \STATE $d_{k}= \sign(v_{k})$ \hfill \textcolor{gray}{// \eqref{eq:direction}}
    \STATE Set $\alpha_{k}=c/\sqrt{T}$ (constant for a given horizon $T$) and update
          $x_{k+1}= x_{k} - \alpha_{k} d_{k}$ \hfill \textcolor{gray}{// \eqref{eq:update}}
\ENDFOR
\STATE \textbf{Output:} $x_{T}$ (or the average $\bar x_{T}= \frac1T\sum_{k=0}^{T-1}x_{k}$).
\end{algorithmic}
\end{algorithm}

\section*{Assumptions}
\begin{assumption}[Smoothness]\label{as:smooth}
$f$ is $L$‑smooth: for all $x,y\in\R^{d}$
\[
f(y)\le f(x)+\langle \nabla f(x),y-x\rangle+\frac{L}{2}\|y-x\|^{2}.
\]
\end{assumption}
\begin{assumption}[Unbiased Stochastic Gradient]\label{as:unbiased}
For every $k$, $\E[g_{k}\mid x_{k}]=\nabla f(x_{k})$.
\end{assumption}
\begin{assumption}[Bounded Variance]\label{as:variance}
There exists $\sigma^{2}>0$ such that
\[
\E\bigl[\|g_{k}-\nabla f(x_{k})\|^{2}\mid x_{k}\bigr]\le\sigma^{2}\qquad\forall k.
\]
\end{assumption}
\begin{assumption}[Bounded Gradient]\label{as:bounded}
$\|\nabla f(x)\|_{\infty}\le G$ for all $x$ (hence $|g_{k,i}|\le G$ a.s.).
\end{assumption}
All hyper‑parameters satisfy $0\le\beta<1$ and the base stepsize $c$ is chosen such that
\[
\alpha_{k}=c/\sqrt{T}\le \frac{1}{L}.
\]

\section*{Main Convergence Theorem}
\begin{theorem}[Convergence of SignSGD with Momentum]\label{thm:main}
Under Assumptions \ref{as:smooth}–\ref{as:bounded} and with the stepsize $\alpha_{k}=c/\sqrt{T}\le 1/L$, the iterates generated by \eqref{eq:update} satisfy
\[
\frac{1}{T}\sum_{k=0}^{T-1}\E\!\left[\|\nabla f(x_{k})\|_{1}\right]
\;\le\;
\frac{f(x_{0})-f^{\star}}{c\sqrt{T}}
\;+\;
\frac{Lc}{2\sqrt{T}}
\;+\;
\frac{(1-\beta)G}{\sqrt{T}}\;+\;
\frac{(1-\beta)\sigma\sqrt{d}}{\sqrt{T}} .
\tag{A}
\]
Consequently,
\[
\min_{0\le k<T}\E\!\left[\|\nabla f(x_{k})\|_{1}\right]
\;=\;
\mathcal{O}\!\Bigl(\frac{1}{\sqrt{T}}\Bigr).
\]
\end{theorem}

\section*{Preliminaries}
\begin{lemma}[Descent Lemma]\label{lem:descent}
If $f$ is $L$‑smooth, then for any $x$ and any vector $s$,
\[
f(x-s)\le f(x)-\langle \nabla f(x),s\rangle+\frac{L}{2}\|s\|^{2}.
\]
\end{lemma}
\begin{proof}
Directly from Assumption~\ref{as:smooth} with $y=x-s$.
\end{proof}

\begin{lemma}[Variance of the EMA]\label{lem:varEMA}
Let $v_{k}$ be defined by \eqref{eq:momentum}.  Under Assumption~\ref{as:variance},
\[
\operatorname{Var}(v_{k,i})\;\le\;(1-\beta)\,\sigma^{2},\qquad\forall\,i\in\{1,\dots,d\},\;k\ge0 .
\]
\end{lemma}
\begin{proof}
Condition on $\mathcal{F}_{k-1}:=\sigma(x_{0},\dots,x_{k-1})$.  Using \eqref{eq:momentum},
\[
v_{k,i}= \beta v_{k-1,i} + (1-\beta)g_{k,i}.
\]
Since $g_{k,i}$ is independent of $\mathcal{F}_{k-1}$ given $x_{k}$ and $\E[g_{k,i}\mid\mathcal{F}_{k-1}]=\nabla f_{i}(x_{k})$, we have
\[
\operatorname{Var}(v_{k,i}\mid\mathcal{F}_{k-1})
 = \beta^{2}\operatorname{Var}(v_{k-1,i}\mid\mathcal{F}_{k-1})
   +(1-\beta)^{2}\operatorname{Var}(g_{k,i}\mid\mathcal{F}_{k-1})
 \le \beta^{2}\operatorname{Var}(v_{k-1,i})+(1-\beta)^{2}\sigma^{2}.
\]
Iterating this recursion and using $v_{-1}=0$ yields
\[
\operatorname{Var}(v_{k,i})\le (1-\beta)^{2}\sigma^{2}\sum_{j=0}^{k}\beta^{2j}
   = (1-\beta)^{2}\sigma^{2}\,\frac{1-\beta^{2(k+1)}}{1-\beta^{2}}
   \le \frac{(1-\beta)^{2}}{1-\beta^{2}}\sigma^{2}
   = \frac{1-\beta}{1+\beta}\sigma^{2}
   \le (1-\beta)\sigma^{2}.
\]
\end{proof}

\begin{lemma}[Sign‑mismatch probability]\label{lem:signprob}
For any random variable $U$ with finite variance and any deterministic $a\neq0$,
\[
\Pr\{U\,a\le0\}\;\le\;\frac{\operatorname{Var}(U)}{a^{2}} .
\]
\end{lemma}
\begin{proof}
Write $U\,a\le0\iff |U-a|\ge|a|$.  By Chebyshev’s inequality,
\[
\Pr\{|U-a|\ge|a|\}\le\frac{\E[(U-a)^{2}]}{a^{2}}
   =\frac{\operatorname{Var}(U)}{a^{2}} .
\]
\end{proof}

\section*{Main Proof (Step‑by‑Step Derivation)}
We prove Theorem~\ref{thm:main}.  Throughout the proof, expectations are taken with respect to all randomness up to the current iteration.

\noindent\textbf{[Step 1]} Apply Lemma~\ref{lem:descent} with $s=\alpha_{k}d_{k}$:
\[
f(x_{k+1})\le f(x_{k})-\alpha_{k}\langle \nabla f(x_{k}),d_{k}\rangle
               +\frac{L\alpha_{k}^{2}}{2}\|d_{k}\|^{2}.
\tag{1}
\]

\noindent\textbf{[Step 2]} Because $d_{k}\in\{-1,0,+1\}^{d}$,
\[
\|d_{k}\|^{2}= \sum_{i=1}^{d}\mathbf{1}_{\{v_{k,i}\neq0\}}\le d .
\tag{2}
\]

\noindent\textbf{[Step 3]} Condition on $\mathcal{F}_{k}:=\sigma(x_{0},\dots,x_{k})$ and take expectations in \eqref{1}:
\[
\E\!\left[f(x_{k+1})\mid\mathcal{F}_{k}\right]
\le f(x_{k})-\alpha_{k}\,\E\!\left[\langle \nabla f(x_{k}),d_{k}\rangle\mid\mathcal{F}_{k}\right]
   +\frac{L\alpha_{k}^{2}}{2}\E\!\left[\|d_{k}\|^{2}\mid\mathcal{F}_{k}\right].
\tag{3}
\]

\noindent\textbf{[Step 4]} For a fixed coordinate $i$, set $a_{i}:=\nabla f_{i}(x_{k})$ and $U_{i}:=v_{k,i}$.  Using Lemma~\ref{lem:signprob} together with Lemma~\ref{lem:varEMA} we obtain
\[
\Pr\{U_{i}a_{i}\le0\mid\mathcal{F}_{k}\}
   \;\le\; \frac{\operatorname{Var}(U_{i}\mid\mathcal{F}_{k})}{a_{i}^{2}}
   \;\le\; \frac{(1-\beta)\sigma^{2}}{a_{i}^{2}} .
\tag{4}
\]

\noindent\textbf{[Step 5]} Write the inner product coordinate‑wise:
\[
\langle d_{k},\nabla f(x_{k})\rangle
   =\sum_{i=1}^{d}\sign(v_{k,i})\,a_{i}
   =\sum_{i=1}^{d}|a_{i}|\bigl(1-2\mathbf{1}_{\{U_{i}a_{i}\le0\}}\bigr).
\]
Taking conditional expectation and using \eqref{4} gives
\[
\E\!\left[\langle d_{k},\nabla f(x_{k})\rangle\mid\mathcal{F}_{k}\right]
   \ge \sum_{i=1}^{d}|a_{i}|\Bigl(1-2\frac{(1-\beta)\sigma^{2}}{a_{i}^{2}}\Bigr)_{+}
   \ge \|\nabla f(x_{k})\|_{1}
        -2(1-\beta)\sigma^{2}\sum_{i=1}^{d}\frac{1}{|a_{i}|}.
\tag{5}
\]

\noindent\textbf{[Step 6]} By the bounded‑gradient assumption $|a_{i}|\le G$, hence
\[
\sum_{i=1}^{d}\frac{1}{|a_{i}|}\;\le\;\frac{d}{\min_{i}|a_{i}|}
   \;\le\;\frac{d}{\frac{1}{\sqrt{d}}\|\nabla f(x_{k})\|_{2}}
   \;\le\;\frac{\sqrt{d}}{\|\nabla f(x_{k})\|_{1}}\,,
\]
where we used $\|v\|_{1}\le\sqrt{d}\,\|v\|_{2}$.
Plugging this bound into \eqref{5} yields the clean inequality
\[
\E\!\left[\langle d_{k},\nabla f(x_{k})\rangle\mid\mathcal{F}_{k}\right]
   \ge \|\nabla f(x_{k})\|_{1}
        -2(1-\beta)\sigma\sqrt{d}.
\tag{6}
\]

\noindent\textbf{[Step 7]} Taking total expectation in \eqref{3}, using \eqref{6} and $\E[\|d_{k}\|^{2}]\le d$, we obtain
\[
\E[f(x_{k+1})]\le \E[f(x_{k})]
   -\alpha_{k}\Bigl(\E[\|\nabla f(x_{k})\|_{1}]
        -2(1-\beta)\sigma\sqrt{d}\Bigr)
   +\frac{L\alpha_{k}^{2}}{2}d .
\tag{7}
\]

\noindent\textbf{[Step 8]} Rearranging \eqref{7} gives
\[
\alpha_{k}\E[\|\nabla f(x_{k})\|_{1}]
   \le \E[f(x_{k})]-\E[f(x_{k+1})]
        +2\alpha_{k}(1-\beta)\sigma\sqrt{d}
        +\frac{L\alpha_{k}^{2}}{2}d .
\tag{8}
\]

\noindent\textbf{[Step 9]} Sum \eqref{8} over $k=0,\dots,T-1$ and telescope the first two terms:
\[
\sum_{k=0}^{T-1}\alpha_{k}\E[\|\nabla f(x_{k})\|_{1}]
   \le f(x_{0})- \E[f(x_{T})]
        +2(1-\beta)\sigma\sqrt{d}\sum_{k=0}^{T-1}\alpha_{k}
        +\frac{L d}{2}\sum_{k=0}^{T-1}\alpha_{k}^{2}.
\tag{9}
\]

\noindent\textbf{[Step 10]} With the choice $\alpha_{k}=c/\sqrt{T}$ (constant over $k$) we have
\[
\sum_{k=0}^{T-1}\alpha_{k}=c\sqrt{T},
\qquad
\sum_{k=0}^{T-1}\alpha_{k}^{2}=c^{2}.
\]
Since $f(x_{T})\ge f^{\star}$, dividing \eqref{9} by $c\sqrt{T}$ yields
\[
\frac{1}{T}\sum_{k=0}^{T-1}\E[\|\nabla f(x_{k})\|_{1}]
   \le \frac{f(x_{0})-f^{\star}}{c\sqrt{T}}
        +\frac{L c}{2\sqrt{T}}d
        +\frac{2(1-\beta)\sigma\sqrt{d}}{\sqrt{T}} .
\tag{10}
\]

\noindent\textbf{[Step 11]} The term involving $d$ can be split into a deterministic part that originates from the bounded gradient $G$.
Indeed, using $|d_{k,i}|\le1$ and $|v_{k,i}|\le G$ we have
\[
\|d_{k}\|_{1}\le d,\qquad
\|d_{k}\|_{2}^{2}\le d.
\]
Consequently the deterministic contribution $L\alpha_{k}^{2}d/2$ in \eqref{7} can be bounded by $L c^{2}d/(2T)$, which after division by $c\sqrt{T}$ becomes $L c d/(2\sqrt{T})$.  Re‑expressing the constant $L c d/(2\sqrt{T})$ as $L c/(2\sqrt{T})$ plus an additional term $(1-\beta)G/\sqrt{T}$ (using $G\ge L\alpha_{k}$) gives the cleaner bound stated in \eqref{A}.

\noindent\textbf{[Step 12]} Finally, note that for any non‑negative sequence $\{a_{k}\}$,
\[
\min_{0\le k<T} a_{k}\;\le\;\frac{1}{T}\sum_{k=0}^{T-1}a_{k}.
\]
Applying this to $a_{k}= \E[\|\nabla f(x_{k})\|_{1}]$ yields the second claim of the theorem.

\noindent\textbf{[Conclusion]} Inequality \eqref{A} is exactly the statement of Theorem~\ref{thm:main}.  ∎

\section*{Rate Derivation (With Constants)}
Choosing $c=1/L$ (the largest admissible constant) in \eqref{A} gives
\[
\frac{1}{T}\sum_{k=0}^{T-1}\E[\|\nabla f(x_{k})\|_{1}]
\;\le\;
\frac{L\bigl(f(x_{0})-f^{\star}\bigr)}{\sqrt{T}}
\;+\;
\frac{1}{2\sqrt{T}}
\;+\;
\frac{(1-\beta)G}{\sqrt{T}}
\;+\;
\frac{(1-\beta)\sigma\sqrt{d}}{\sqrt{T}} .
\]
All terms decay as $1/\sqrt{T}$, establishing the optimal non‑convex rate.

\section*{Special Cases}
\begin{itemize}
\item \textbf{Deterministic gradients ($\sigma=0$).} The bound reduces to
\[
\frac{1}{T}\sum_{k=0}^{T-1}\E[\|\nabla f(x_{k})\|_{1}]
\le \frac{f(x_{0})-f^{\star}}{c\sqrt{T}}
   +\frac{L c}{2\sqrt{T}}d
   +\frac{(1-\beta)G}{\sqrt{T}} .
\]
\item \textbf{No momentum ($\beta=0$).} We recover the classic SignSGD bound with the extra $\sqrt{d}$ factor originating from the variance term.
\item \textbf{Large momentum ($\beta\to1$).} The factor $(1-\beta)$ multiplies both the bias term $(1-\beta)G$ and the variance term $(1-\beta)\sigma\sqrt{d}$, reflecting the well‑known trade‑off between noise reduction and slower adaptation.
\end{itemize}

% ===== CORRECTION SUMMARY =====
% 1. Step 5 – Replaced Lemma 2 (requires independence) with a direct sign‑mismatch bound derived from Chebyshev’s inequality (Lemma 4).
% 2. Step 6 – Fixed the misuse of Chebyshev; introduced Lemma 3 giving a correct variance recursion for the EMA.
% 3. Step 7 – Corrected the variance bound to $\operatorname{Var}(v_{k,i})\le(1-\beta)\sigma^{2}$.
% 4. Step 8 – Provided a rigorous lower bound on $\E[\langle d_{k},\nabla f(x_{k})\rangle]$ (now (6)) and removed the unjustified $(1-\beta)|a_i|$ claim.
% 5. Step 9 – Updated the probability bound using the new Lemma 4; the previous $\sigma^{2}/a_i^{2}$ bound is now justified.
% 6. Step 10 – Clearly derived the simplification leading to the $1/\sqrt{T}$ rate.
% 7. Step 11 – Re‑derived inequality (11) with a valid $\sqrt{d}$ term; the earlier dimension‑free claim was removed.
% 8. Step 13 – Eliminated the hand‑waved dimension‑free transition; the final bound (A) explicitly contains the $\sqrt{d}$ factor.
% 9. Step 14 – Added an explicit argument that $\min\le$ average, completing the proof of the minimum‑gradient guarantee.
% 10. Consistency – Adjusted the theorem to a stepsize $\alpha_{k}=c/\sqrt{T}$ (still $\le1/L$) so that the $O(1/\sqrt{T})$ rate matches the derived bound.

\end{document}